\PassOptionsToPackage{table}{xcolor}
\documentclass[a4paper]{article}

\usepackage[margin=1.8cm]{geometry}
\usepackage{graphicx}
\usepackage{amsmath,amssymb}
\usepackage{minted}
\usepackage{multirow}
\usepackage{float}
\usepackage{todonotes}
\usepackage{forest}
\usepackage{xstring}
\usepackage{xcolor}
\usepackage[most]{tcolorbox}
\usepackage{xparse}
\usepackage{natbib}

% Packages used for responsive tables and text wrapping
\usepackage{booktabs}
\usepackage{tabularx}
\usepackage{array}
\usepackage{ragged2e}
\usepackage{microtype}
\usepackage[hidelinks]{hyperref}

\usetikzlibrary{graphs,graphdrawing}
\usegdlibrary{layered}

% for external graphviz
\input{/home/donkarlo/Dropbox/repo/nd_ai_project/src/nd_ai/knoweldge_representation/language/formal/computer/programming/tex/latex/code_clips/external_graphviz.tex}

% for tags
\input{/home/donkarlo/Dropbox/repo/nd_ai_project/src/nd_ai/knoweldge_representation/language/formal/computer/programming/tex/latex/parts/control_sequence/kind/semtag/semtag.tex}

% for links
\input{/home/donkarlo/Dropbox/repo/nd_ai_project/src/nd_ai/knoweldge_representation/language/formal/computer/programming/tex/latex/code_clips/seehere.tex}

% ------------------------------------------------
% Document style
% ------------------------------------------------

\definecolor{headerblue}{HTML}{1F3A5F}
\definecolor{rowgray}{HTML}{F3F5F7}
\definecolor{lightblue}{HTML}{EEF4FA}

\setlength{\parindent}{0pt}
\setlength{\parskip}{6pt}
\setlength{\emergencystretch}{3em}

\renewcommand{\arraystretch}{1.35}

\newcolumntype{Y}{>{\RaggedRight\arraybackslash}X}
\newcolumntype{T}{>{\RaggedRight\arraybackslash}p{3.2cm}}

\title{
    \textbf{FastAPI and Django}\[4pt]
        {\large Their Purpose, Differences, and Relationship}
}
\date{}
\author{Mohammad Rahmani}

\begin{document}
    
    \maketitle
    
    
    \section*{What Is FastAPI?}
        
        FastAPI is an independent Python web framework designed primarily for
        building web APIs.
        
        An API allows another program, frontend application, mobile application,
        or machine-learning service to communicate with the backend through
        HTTP requests.
        
        FastAPI helps a developer define:
        
        \begin{itemize}
            \item API addresses such as \texttt{/users} or \texttt{/predict};
            \item accepted input parameters;
            \item expected data types;
            \item validation rules;
            \item JSON responses;
            \item error responses;
            \item automatic API documentation.
        \end{itemize}
        
        FastAPI uses normal Python type hints to understand the expected input
        and output data. It can automatically validate incoming data and reject
        invalid requests before they reach the main application logic.
        
        \begin{tcolorbox}[
            colback=lightblue,
            colframe=headerblue,
            title=\textbf{Important},
            fonttitle=\bfseries,
            breakable
        ]
            FastAPI is not a deployment tool like Docker.
            
            FastAPI is used to write the API application itself. Docker can then be
            used to package and run that FastAPI application in a reproducible
            container.
        \end{tcolorbox}
    
    
    \section*{A Simple FastAPI Example}
        
        \begin{minted}[
            fontsize=\small,
            breaklines,
            breakanywhere,
            frame=single,
            framesep=3mm
        ]{python}
from fastapi import FastAPI

app = FastAPI()

@app.get("/hello")
def hello():
return {"message": "Hello from FastAPI"}
        \end{minted}
        
        When a client sends a GET request to \texttt{/hello}, FastAPI returns:
        
        \begin{minted}[
            fontsize=\small,
            breaklines,
            breakanywhere,
            frame=single,
            framesep=3mm
        ]{json}
{
"message": "Hello from FastAPI"
}
        \end{minted}
    
    
    \section*{What Does FastAPI Provide?}
        
        \rowcolors{2}{rowgray}{white}
        
        \begin{tabularx}{\textwidth}{T Y}

\toprule
\rowcolor{headerblue}
\color{white}\textbf{Feature} &
\color{white}\textbf{Explanation} \
\midrule

Routing &
Connects addresses such as \texttt{/users}, \texttt{/orders}, or
\texttt{/predict} to Python functions. \

Input validation &
Checks whether request data has the correct structure and data types. \

JSON serialization &
Converts Python objects into JSON responses that other applications can
understand. \

Type-hint integration &
Uses Python type annotations to define parameters, request bodies, and
response structures. \

Automatic documentation &
Automatically produces interactive OpenAPI documentation using Swagger
UI and ReDoc. \

Asynchronous endpoints &
Supports \texttt{async} functions for operations involving databases,
networks, files, and external services. \

Dependency injection &
Provides a structured way to supply database sessions, authenticated
users, configuration, and reusable logic to endpoints. \

Error handling &
Converts application and validation errors into structured HTTP
responses. \

\bottomrule
        
        \end{tabularx}
    
    
    \section*{Does Django Use FastAPI?}
        
        \begin{tcolorbox}[
            colback=rowgray,
            colframe=headerblue,
            boxrule=0.7pt,
            arc=2pt,
            breakable
        ]
            \centering
            \Large\textbf{No. Django does not use FastAPI.}
        \end{tcolorbox}
        
        FastAPI and Django are two separate Python web frameworks.
        
        A normal Django application does not run FastAPI internally, and FastAPI
        does not require Django. A developer can choose either framework based
        on the requirements of the project.
        
        It is technically possible to place Django and FastAPI in the same large
        system as separate services, but this is an architectural decision. It
        does not mean that Django is based on FastAPI.
    
    
    \section*{What Does Django Use Instead?}
        
        Django has its own web framework components, including:
        
        \begin{itemize}
            \item URL routing;
            \item function-based and class-based views;
            \item HTTP request and response objects;
            \item JSON responses;
            \item database models and an ORM;
            \item authentication and authorization;
            \item an administration panel;
            \item HTML templates;
            \item forms and input validation;
            \item middleware.
        \end{itemize}
        
        For a simple JSON API, Django can return a \texttt{JsonResponse} directly.
    
    
    \section*{A Simple Django API Example}
        
        \begin{minted}[
            fontsize=\small,
            breaklines,
            breakanywhere,
            frame=single,
            framesep=3mm
        ]{python}
from django.http import JsonResponse

def hello(request):
return JsonResponse(
{"message": "Hello from Django"}
)
        \end{minted}
        
        The corresponding Django URL configuration could be:
        
        \begin{minted}[
            fontsize=\small,
            breaklines,
            breakanywhere,
            frame=single,
            framesep=3mm
        ]{python}
from django.urls import path
from .views import hello

urlpatterns = [
path("hello/", hello),
]
        \end{minted}
    
    
    \section*{Django REST Framework}
        
        Django itself can return JSON responses, but larger REST APIs usually
        need additional features such as:
        
        \begin{itemize}
            \item serializers;
            \item API authentication;
            \item permissions;
            \item pagination;
            \item filtering;
            \item browsable API pages;
            \item view sets;
            \item structured validation.
        \end{itemize}
        
        A common solution is \textbf{Django REST Framework}, often abbreviated
        as \textbf{DRF}.
        
        DRF is installed on top of Django and uses Django's models, views,
        authentication system, and request-response infrastructure.
        
        \begin{tcolorbox}[
            colback=lightblue,
            colframe=headerblue,
            title=\textbf{Do not confuse these two statements},
            fonttitle=\bfseries,
            breakable
        ]
            \textbf{Django has its own basic API capabilities.}
            
            However, \textbf{Django REST Framework is a separate package} that adds
            more specialized tools for building REST APIs.
        \end{tcolorbox}
    
    
    \section*{FastAPI Versus Django}
        
        \begin{tabularx}{\textwidth}{T Y Y}

\toprule
\rowcolor{headerblue}
\color{white}\textbf{Aspect} &
\color{white}\textbf{FastAPI} &
\color{white}\textbf{Django} \
\midrule

Primary purpose &
Building APIs and backend services. &
Building complete web applications and backend systems. \

Relationship &
Independent Python framework. &
Independent Python framework; it does not use FastAPI. \

API support &
API functionality is the central purpose of the framework. &
Can create APIs directly or through Django REST Framework. \

Database ORM &
Does not include a full built-in ORM; it is commonly combined with
SQLAlchemy or SQLModel. &
Includes Django ORM as a central part of the framework. \

Administration panel &
Does not provide a Django-style administration interface by default. &
Includes an automatically generated administration panel. \

HTML templates &
Possible, but not its primary focus. &
Includes a complete template system for HTML websites. \

Authentication &
Provides tools and patterns, but the developer usually assembles the
required authentication system. &
Includes users, sessions, permissions, password handling, and
authentication components. \

Automatic API documentation &
Built in through OpenAPI, Swagger UI, and ReDoc. &
Not equivalent in basic Django; API documentation usually requires
additional packages or DRF tooling. \

Typical use &
Microservices, ML model serving, mobile backends, and standalone APIs. &
Business applications, content systems, dashboards, administrative
systems, and full websites. \

\bottomrule
        
        \end{tabularx}
    
    
    \section*{Which One Is More Suitable?}
        
        \subsection*{Choose FastAPI When}
            
            \begin{itemize}
                \item the main product is an API;
                \item you are building a microservice;
                \item you want to expose a machine-learning model;
                \item automatic OpenAPI documentation is important;
                \item the frontend and backend are separate applications;
                \item you want a relatively small and focused backend framework.
            \end{itemize}
        
        \subsection*{Choose Django When}
            
            \begin{itemize}
                \item you need a complete website and backend;
                \item you need a built-in administration panel;
                \item the application has many related database models;
                \item you need users, sessions, permissions, and forms;
                \item you want many web-development features in one framework;
                \item the project is a large business or content-management system.
            \end{itemize}
        
        \subsection*{Choose Django REST Framework When}
            
            \begin{itemize}
                \item the project already uses Django;
                \item you want to expose Django models through a REST API;
                \item you need serializers, permissions, pagination, and view sets;
                \item you want Django's complete ecosystem together with a mature
                REST API toolkit.
            \end{itemize}
    
    
    \section*{Relationship with Docker}
        
        FastAPI, Django, and Docker operate at different levels:
        
        \begin{tcolorbox}[
            colback=rowgray,
            colframe=headerblue,
            boxrule=0.7pt,
            arc=2pt,
            breakable
        ]
            \begin{center}
                \textbf{FastAPI or Django}\
                Application framework used to write the backend
            \end{center}
            
            [
            \Downarrow
            ]
            
            \begin{center}
                \textbf{Docker}\
                Packages the application, Python runtime, libraries, and configuration
            \end{center}
            
            [
            \Downarrow
            ]
            
            \begin{center}
                \textbf{Server, cloud platform, or Kubernetes}\
                Runs and manages the packaged application
            \end{center}
        \end{tcolorbox}
    
    
    \section*{Final Summary}
        
        \begin{tcolorbox}[
            colback=lightblue,
            colframe=headerblue,
            title=\textbf{Summary},
            fonttitle=\bfseries,
            breakable
        ]
            FastAPI is a Python framework specialized in building APIs.
            
            Django does not use FastAPI. Django has its own routing, views, request
            and response classes, ORM, authentication system, templates, and
            administration panel.
            
            Django can create simple JSON APIs using \texttt{JsonResponse}. For more
            advanced REST APIs, Django projects commonly use the separate
            \textbf{Django REST Framework} package.
            
            FastAPI is usually a suitable choice for focused APIs, microservices,
            and machine-learning model serving. Django is usually a suitable choice
            for complete web applications that need database models, users,
            administration, forms, and HTML pages.
        \end{tcolorbox}
        
%        \bibliographystyle{plainnat}
%        \bibliography{/home/donkarlo/Dropbox/repo/nd_shared_research/bibliography/bibliography.bib}

\end{document}

